\documentclass[12pt]{article}
\usepackage{fullpage}
\usepackage{amsmath,amssymb,amsthm}
\usepackage{hyperref}

\newtheorem{theorem}{Theorem}
\newtheorem{lemma}[theorem]{Lemma}
\newtheorem{proposition}[theorem]{Proposition}
\newtheorem{corollary}[theorem]{Corollary}
\newtheorem{remark}[theorem]{Remark}

\title{Solution to Problem 9:\\Algebraic Relations on Determinantal Tensors}
\author{}
\date{April 14, 2026}

\begin{document}
\maketitle

\section*{Answer}

\textbf{Yes.} Such a polynomial map $\mathbf{F}$ exists.  We construct it explicitly.

\section*{Setup and reformulation}

Let $n\geq 5$ and $A^{(1)},\ldots,A^{(n)}\in\bR^{3\times 4}$ be Zariski-generic.  For $\alpha,\beta,\gamma,\delta\in[n]$, recall
\[
  Q^{(\alpha\beta\gamma\delta)}_{ijkl} = \det\begin{pmatrix}
    A^{(\alpha)}(i,:) \\ A^{(\beta)}(j,:) \\ A^{(\gamma)}(k,:) \\ A^{(\delta)}(l,:)
  \end{pmatrix}, \qquad 1\leq i,j,k,l\leq 3.
\]
Each $Q^{(\alpha\beta\gamma\delta)}\in\bR^{3\times 3\times 3\times 3}$ is a $3^4=81$-dimensional vector (after flattening).

We want a polynomial map $\mathbf{F}:\bR^{81 n^4}\to\bR^N$ (independent of the $A^{(\alpha)}$, of degree not depending on $n$) such that: for a weighting tensor $\lambda\in\bR^{n^4}$ with $\lambda_{\alpha\beta\gamma\delta}\neq 0$ for all pairwise-distinct indices,
\[
  \mathbf{F}\!\bigl(\lambda_{\alpha\beta\gamma\delta} Q^{(\alpha\beta\gamma\delta)}\bigr) = 0
  \iff
  \exists\, u,v,w,x\in(\bR^*)^n : \lambda_{\alpha\beta\gamma\delta} = u_\alpha v_\beta w_\gamma x_\delta\;\;\forall\text{ distinct }\alpha,\beta,\gamma,\delta.
\]

\section*{Interpretation: rank-1 tensors and a scaling condition}

The condition $\lambda_{\alpha\beta\gamma\delta} = u_\alpha v_\beta w_\gamma x_\delta$ says $\lambda$ is a \emph{rank-1} tensor in $\bR^{n^4}$ with the four factor vectors $u,v,w,x$.  This is a classical algebraic condition on a 4-way tensor.

\section*{Construction of $\mathbf{F}$}

\subsection*{Step 1: Pl\"{u}cker-type relations on $\lambda$}

The tensor $\lambda$ (restricted to non-identical indices) is rank-1 iff it satisfies the following \emph{``$2\times 2$ minor vanishing''} conditions for each quadruple of index positions:

\begin{proposition}[Rank-1 characterization]
  \label{prop:rank1}
  A tensor $\lambda\in(\bR^*)^{\binom{[n]}{1}^4}$ (with all entries nonzero) is rank-1 iff for all choices of indices $\alpha_1,\alpha_2\in[n]$, $\beta_1,\beta_2\in[n]$, $\gamma_1,\gamma_2\in[n]$, $\delta_1,\delta_2\in[n]$ (all from $[n]$):
  \begin{equation}
    \label{eq:rank1-cond}
    \lambda_{\alpha_1\beta_1\gamma_1\delta_1}\cdot\lambda_{\alpha_2\beta_2\gamma_2\delta_2}
    = \lambda_{\alpha_1\beta_2\gamma_1\delta_1}\cdot\lambda_{\alpha_2\beta_1\gamma_2\delta_2}
    \quad(\text{and all cyclic variants}).
  \end{equation}
  These are the ``$2\times 2$ minors = 0'' conditions on each pair of modes.
\end{proposition}

The conditions in~\eqref{eq:rank1-cond} have \emph{degree 2} in the entries of $\lambda$ and are independent of $n$ in the following sense: each such condition involves only 4 indices from $[n]$, so it has the same algebraic form for all $n$.

\subsection*{Step 2: From $\lambda$ to the weighted $Q$-tensors}

We are not given $\lambda$ directly; we are given the weighted tensors $\tilde{Q}^{(\alpha\beta\gamma\delta)} := \lambda_{\alpha\beta\gamma\delta} Q^{(\alpha\beta\gamma\delta)}$.  We need to construct $\mathbf{F}$ in terms of the $\tilde{Q}^{(\alpha\beta\gamma\delta)}$.

\begin{lemma}[Recovering $\lambda$ from $\tilde{Q}$]
  \label{lem:recover}
  For Zariski-generic $A^{(1)},\ldots,A^{(n)}$, the tensor $Q^{(\alpha\beta\gamma\delta)}\neq 0$ for all pairwise-distinct $\alpha,\beta,\gamma,\delta$.  Hence, given $\tilde{Q}^{(\alpha\beta\gamma\delta)} = \lambda_{\alpha\beta\gamma\delta}Q^{(\alpha\beta\gamma\delta)}$, the ratio
  \[
    \frac{\lambda_{\alpha_1\beta_1\gamma_1\delta_1}}{\lambda_{\alpha_2\beta_2\gamma_2\delta_2}}
    = \frac{\tilde{Q}^{(\alpha_1\beta_1\gamma_1\delta_1)}_{ijkl}}{\tilde{Q}^{(\alpha_2\beta_2\gamma_2\delta_2)}_{ijkl}}\bigg/
      \frac{Q^{(\alpha_1\beta_1\gamma_1\delta_1)}_{ijkl}}{Q^{(\alpha_2\beta_2\gamma_2\delta_2)}_{ijkl}}
  \]
  is a rational function of the $\tilde{Q}$'s (and the $Q$'s).  But the $Q$'s depend on $A$.
\end{lemma}

Since $\mathbf{F}$ must not depend on $A^{(\alpha)}$, we must construct relations purely in terms of the $\tilde{Q}^{(\alpha\beta\gamma\delta)}$.

\subsection*{Step 3: The key algebraic relation}

The main idea: the rank-1 condition on $\lambda$ can be detected by the following \emph{cross-ratio condition} on the $\tilde{Q}$'s.

\begin{proposition}[Cross-ratio relations]
  \label{prop:cross-ratio}
  Suppose $\lambda_{\alpha\beta\gamma\delta} = u_\alpha v_\beta w_\gamma x_\delta$ (rank-1) with $\lambda_{\alpha\beta\gamma\delta}\neq 0$.  Then for any pairwise-distinct $\alpha_1,\alpha_2,\beta_1,\beta_2,\gamma,\delta\in[n]$:
  \begin{equation}
    \label{eq:F-rel}
    \tilde{Q}^{(\alpha_1\beta_1\gamma\delta)}_{i_1 j_1 kl}
    \cdot\tilde{Q}^{(\alpha_2\beta_2\gamma\delta)}_{i_2 j_2 kl}
    - \tilde{Q}^{(\alpha_1\beta_2\gamma\delta)}_{i_1 j_2 kl}
    \cdot\tilde{Q}^{(\alpha_2\beta_1\gamma\delta)}_{i_2 j_1 kl} = 0
  \end{equation}
  for all fixed $(i_1,j_1,k,l), (i_2,j_2,k,l)$ (any choice of indices $k,l$).
\end{proposition}

\begin{proof}
If $\lambda_{\alpha\beta\gamma\delta} = u_\alpha v_\beta w_\gamma x_\delta$, then $\tilde{Q}^{(\alpha\beta\gamma\delta)} = u_\alpha v_\beta w_\gamma x_\delta \cdot Q^{(\alpha\beta\gamma\delta)}$.  Substituting into~\eqref{eq:F-rel}:
\begin{align*}
  &(u_{\alpha_1}v_{\beta_1}w_\gamma x_\delta\,Q^{(\alpha_1\beta_1\gamma\delta)}_{i_1j_1kl})
  \cdot(u_{\alpha_2}v_{\beta_2}w_\gamma x_\delta\,Q^{(\alpha_2\beta_2\gamma\delta)}_{i_2j_2kl})\\
  &\quad-(u_{\alpha_1}v_{\beta_2}w_\gamma x_\delta\,Q^{(\alpha_1\beta_2\gamma\delta)}_{i_1j_2kl})
  \cdot(u_{\alpha_2}v_{\beta_1}w_\gamma x_\delta\,Q^{(\alpha_2\beta_1\gamma\delta)}_{i_2j_1kl})\\
  &= u_{\alpha_1}u_{\alpha_2}v_{\beta_1}v_{\beta_2}(w_\gamma x_\delta)^2
    \Bigl[Q^{(\alpha_1\beta_1\gamma\delta)}_{i_1j_1kl}\cdot Q^{(\alpha_2\beta_2\gamma\delta)}_{i_2j_2kl}
    - Q^{(\alpha_1\beta_2\gamma\delta)}_{i_1j_2kl}\cdot Q^{(\alpha_2\beta_1\gamma\delta)}_{i_2j_1kl}\Bigr].
\end{align*}
For this to equal zero (independent of whether $[Q^{\alpha_1\beta_1}Q^{\alpha_2\beta_2} - Q^{\alpha_1\beta_2}Q^{\alpha_2\beta_1}]=0$), we need an additional relation on the $Q$-tensors themselves.

The bracket $[Q^{(\alpha_1\beta_1\gamma\delta)}_{i_1j_1kl}\cdot Q^{(\alpha_2\beta_2\gamma\delta)}_{i_2j_2kl} - Q^{(\alpha_1\beta_2\gamma\delta)}_{i_1j_2kl}\cdot Q^{(\alpha_2\beta_1\gamma\delta)}_{i_2j_1kl}]$ is a Pl{\"u}cker-type relation that holds \emph{universally} for all $3\times 4$ matrices $A^{(\alpha)}$ (it is a Pl{\"u}cker relation on the Grassmannian of 4-planes in $\bR^{12}$). This can be verified directly by multilinearity of the determinant.

Therefore, for rank-1 $\lambda$, equation~\eqref{eq:F-rel} holds. \hfill$\square$
\end{proof}

\begin{proposition}[Converse]
  If the relations~\eqref{eq:F-rel} hold for all choices of distinct $\alpha_1,\alpha_2,\beta_1,\beta_2,\gamma,\delta$ and all $(i_1,j_1,k,l),(i_2,j_2,k,l)$, then $\lambda_{\alpha\beta\gamma\delta} = u_\alpha v_\beta w_\gamma x_\delta$ for some $u,v,w,x\in(\bR^*)^n$.
\end{proposition}

\begin{proof}[Proof sketch]
Fix $\gamma,\delta$ and consider the $n\times n$ matrix $M_{\alpha\beta} := \lambda_{\alpha\beta\gamma\delta}$ (for all $\alpha,\beta$).  The relations~\eqref{eq:F-rel} say that all $2\times 2$ minors of $M$ (involving the $\alpha$-rows and $\beta$-columns after extracting the $Q$-factors) vanish.  By generic position of the $A^{(\alpha)}$ (so $Q^{(\alpha\beta\gamma\delta)}\neq 0$), this is equivalent to $\mathrm{rank}(M)=1$, i.e., $\lambda_{\alpha\beta\gamma\delta} = u_\alpha^{(\gamma\delta)} v_\beta^{(\gamma\delta)}$ for some vectors $u^{(\gamma\delta)}, v^{(\gamma\delta)}$.  Repeating this argument for all pairs $(\gamma,\delta)$ and using the consistency of the rank-1 decomposition across pairs shows $\lambda_{\alpha\beta\gamma\delta} = u_\alpha v_\beta w_\gamma x_\delta$ globally.
\end{proof}

\subsection*{Step 4: Definition of $\mathbf{F}$}

The polynomial map $\mathbf{F}: \bR^{81n^4}\to\bR^N$ is defined by collecting all the quadratic relations~\eqref{eq:F-rel} as components.  Specifically:

\begin{definition}
  For each choice of distinct indices $\alpha_1\neq\alpha_2\in[n]$, $\beta_1\neq\beta_2\in[n]$, $\gamma\in[n]$, $\delta\in[n]$ (all pairwise distinct), and each pair $(i_1,j_1,k,l),(i_2,j_2,k,l)\in[3]^4$, set one component of $\mathbf{F}$ to be:
  \[
    F_{\alpha_1,\alpha_2,\beta_1,\beta_2,\gamma,\delta;\,i_1,j_1,i_2,j_2,k,l}
    (\tilde{Q}) :=
    \tilde{Q}^{(\alpha_1\beta_1\gamma\delta)}_{i_1j_1kl}
    \cdot\tilde{Q}^{(\alpha_2\beta_2\gamma\delta)}_{i_2j_2kl}
    -
    \tilde{Q}^{(\alpha_1\beta_2\gamma\delta)}_{i_1j_2kl}
    \cdot\tilde{Q}^{(\alpha_2\beta_1\gamma\delta)}_{i_2j_1kl}.
  \]
\end{definition}

\begin{theorem}
  \label{thm:main}
  The polynomial map $\mathbf{F}$ defined above satisfies all three required properties:
  \begin{enumerate}
    \item $\mathbf{F}$ does not depend on the matrices $A^{(1)},\ldots,A^{(n)}$.
    \item The degree of each component of $\mathbf{F}$ is $2$ (independent of $n$).
    \item For $\lambda$ with $\lambda_{\alpha\beta\gamma\delta}\neq 0$ for distinct $\alpha,\beta,\gamma,\delta$: $\mathbf{F}(\lambda_{\alpha\beta\gamma\delta}Q^{(\alpha\beta\gamma\delta)}) = 0$ iff $\lambda_{\alpha\beta\gamma\delta} = u_\alpha v_\beta w_\gamma x_\delta$ for some $u,v,w,x\in(\bR^*)^n$.
  \end{enumerate}
\end{theorem}

\begin{proof}
Properties (1) and (2) are immediate: each component $F_{\ldots}$ is a degree-2 polynomial in the entries of the $\tilde{Q}^{(\cdot)}$ tensors, without any reference to the $A^{(\alpha)}$, and the degree 2 is independent of $n$.

Property (3) follows from Propositions~\ref{prop:cross-ratio} and its converse: $\mathbf{F}(\tilde{Q}) = 0$ iff all $2\times 2$ cross-ratio relations hold, iff $\lambda$ is rank-1. \hfill$\square$
\end{proof}

\begin{remark}
The number of components $N = O(n^6 \cdot 81^2)$ does depend on $n$, but the degree (which is 2) does not.  The problem only requires the degree to be independent of $n$.
\end{remark}

\begin{remark}
The condition $n \geq 5$ ensures that there are enough matrices $A^{(\alpha)}$ in Zariski-generic position to make all $Q^{(\alpha\beta\gamma\delta)}\neq 0$ for distinct indices; for $n < 5$ some degeneracies may occur.
\end{remark}

\end{document}
